\section*{Ethics Statement}

\paragraph{No evaluation with Deaf users.} We did not evaluate the system with Deaf or
hard-of-hearing users, and no Deaf signer took part in its design. Sign language technology
should be built with, and ideally led by, Deaf communities \citep{bragg2019interdisciplinary,
yin2021including,fox2023bestpractices}, and Deaf researchers have criticized how often the field
sets its agenda without them \citep{desai2024systemic}. This paper is an engineering study of one
failure mode, the train--serve gap, and not a claim of usefulness to the people the technology is
meant for.

\paragraph{Practice mode can mislead learners.} Deaf children of hearing parents are at risk of
language deprivation when no one around them signs \citep{humphries2012deprivation}, and practice
tools are one motivation for this work. At the live accuracy we measured, however, the Practice
mode would tell a learner that many correctly produced signs are wrong, and for some signs it fails
every time; it could discourage learners or teach them to distort signs to please the model. It
should report low-confidence results as ``not sure'' rather than ``try again'', state per-sign
reliability, and be evaluated with learners and Deaf teachers before anyone relies on it.

\paragraph{Overstating translation.} Calling a gloss classifier a translator misleads users about
what it understands, and the app's main mode is still named \emph{Translate}. The app states that it
recognizes a fixed set of isolated signs and is not a full ASL translator, shows ``not sure'' below
its threshold, and labels rule-based sentences as basic, but a fluent generated sentence can lend
false authority to wrong glosses. A tool that misrecognizes signs could cause harm if relied on in
important conversations; it must not replace interpreters, and a name such as ``Recognize'' would
describe the mode better.

\paragraph{Data and consent.} The training data is the public Google ISLR release (CC BY 4.0) of
landmarks, not video. The live clips show the author of this paper, who recorded them and consents
to their use here; they are not released, and every figure in the paper shows landmarks only,
never camera pixels. The app processes video on the device; when the user asks for a sentence, the
recognized glosses, and nothing else, are sent to a third-party language model provider.
